% Standalone problem document, generated from the adjacent README.md.
% Compile with XeLaTeX. Mathematical content is maintained in Markdown.
\documentclass[11pt,a4paper]{article}
\usepackage[fontsize=10.5pt]{scrextend}
\usepackage[margin=22mm,headheight=15pt,headsep=9mm,footskip=11mm]{geometry}
\usepackage{fontspec}
\setmainfont{texgyrepagella}[Extension=.otf,UprightFont=*-regular,BoldFont=*-bold,ItalicFont=*-italic,BoldItalicFont=*-bolditalic]
\setsansfont{texgyreheros}[Extension=.otf,UprightFont=*-regular,BoldFont=*-bold,ItalicFont=*-italic,BoldItalicFont=*-bolditalic]
\setmonofont{texgyrecursor}[Extension=.otf,UprightFont=*-regular,BoldFont=*-bold,ItalicFont=*-italic,BoldItalicFont=*-bolditalic,Scale=MatchLowercase]
\usepackage{amsmath,amssymb}
\usepackage{unicode-math}
\setmathfont{texgyrepagella-math.otf}
\usepackage{xcolor,microtype,enumitem,needspace,fancyhdr}
\usepackage{bookmark}
\definecolor{ink}{HTML}{17344B}
\definecolor{accent}{HTML}{197D80}
\definecolor{muted}{HTML}{596773}
\hypersetup{colorlinks=true,linkcolor=accent,urlcolor=accent,pdftitle={MF-11 - Temporal
regularity of marginal matrix-product
growth},pdfauthor={Open Problems in Numerical Linear Algebra}}
\urlstyle{same}
\setlength{\parindent}{0pt}
\setlength{\parskip}{5pt plus 1pt minus 1pt}
\linespread{1.04}
\setlength{\emergencystretch}{3em}
\widowpenalty=10000
\clubpenalty=10000
\displaywidowpenalty=10000
\allowdisplaybreaks[1]
\setlist{nosep,leftmargin=1.5em}
\providecommand{\tightlist}{\setlength{\itemsep}{0pt}\setlength{\parskip}{0pt}}
\providecommand{\pandocbounded}[1]{#1}
\pagestyle{fancy}
\fancyhf{}
\fancyhead[L]{\sffamily\footnotesize\color{muted}OPEN PROBLEMS IN NUMERICAL LINEAR ALGEBRA}
\fancyhead[R]{\sffamily\bfseries\footnotesize\color{accent}MF-11}
\fancyfoot[L]{\parbox[b]{0.9\textwidth}{\sffamily\fontsize{8}{10}\selectfont\color{muted}Editorial ratings; a bounded search cannot certify that a problem remains unsolved.\\Literature check: 2026-09-10}}
\fancyfoot[R]{\sffamily\footnotesize\color{muted}\thepage}
\renewcommand{\headrulewidth}{0.3pt}
\renewcommand{\headrule}{\hbox to\headwidth{\color{accent}\leaders\hrule height \headrulewidth\hfill}}
\makeatletter
\renewcommand{\section}{\Needspace{5\baselineskip}\@startsection{section}{1}{0pt}{12pt plus 2pt minus 2pt}{4pt}{\sffamily\bfseries\large\color{ink}}}
\renewcommand{\subsection}{\Needspace{4\baselineskip}\@startsection{subsection}{2}{0pt}{9pt plus 1pt minus 1pt}{3pt}{\sffamily\bfseries\normalsize\color{ink}}}
\makeatother
\setcounter{secnumdepth}{0}
\begin{document}
{\sffamily\small\color{muted}Matrix functions and stability\par}
\vspace{3pt}
{\raggedright\hyphenpenalty=10000\exhyphenpenalty=10000\sffamily\bfseries\fontsize{21}{24}\selectfont\color{ink}Temporal
regularity of marginal matrix-product growth\par}
\vspace{7pt}
{\sffamily\small\textbf{Difficulty:} challenging\quad\textbf{Importance:} interesting
to specialist\par}
{\sffamily\small\bfseries\color{accent}Status: Open\par}
\vspace{4pt}
{\color{accent}\hrule height 0.8pt}
\vspace{6pt}
\textbf{Rating rationale:} Challenging because cancellations across
coupled blocks resist general growth regularity; its immediate impact is
on specialists in marginally stable matrix products.

\subsection{Context and notation}\label{context-and-notation}

The joint spectral radius of a nonempty compact set
\(\mathcal M\subset\mathbb C^{d\times d}\) is

\[
\widehat\rho(\mathcal M)=\lim_{k\to\infty}
\max_{A_1,\ldots,A_k\in\mathcal M}\|A_k\cdots A_1\|_2^{1/k}.
\]

This definition also applies to finite real matrix sets. The ordinary
spectral radius of one matrix is written \(\rho(A)\).

For a compact nonempty \(\mathcal M\subset\mathbb R^{d\times d}\) with
\(\widehat\rho(\mathcal M)=1\), define its maximal product norm at
length \(k\) by

\[
g_{\mathcal M}(k)=\max_{A_1,\ldots,A_k\in\mathcal M}\|A_k\cdots A_1\|_2.
\]

This problem concerns growth within one family over time.
\href{https://github.com/ajt60gaibb/OpenProblemsInNLA/blob/main/matrix-functions-and-stability/MF-07/README.md}{MF-07}
instead requests a dimension-dependent bound uniform across families.

\subsection{Problem statement}\label{problem-statement}

Does every family in this notation satisfy both

\[
\exists c>0\ \forall k,\ell\ge1:\quad
g_{\mathcal M}(k+\ell)\ge c\,g_{\mathcal M}(k)
\]

and

\[
\forall \ell\ge1\ \exists C_\ell>0\ \forall k\ge1:\quad
g_{\mathcal M}(\ell k)\le C_\ell\,g_{\mathcal M}(k)?
\]

Constants may depend on \(\mathcal M\).

\subsection{Reference and status
evidence}\label{reference-and-status-evidence}

Varney and Morris, \href{https://arxiv.org/html/2209.00449}{On marginal
growth rates of matrix products}, Linear Algebra Appl. 709 (2025),
132--163, \href{https://doi.org/10.1016/j.laa.2025.01.013}{DOI}, §4,
Definition 4.1 and §7, Question 1. The source names these properties
weak increase and weak upper regular variation. Searches for both
property names, the authors, and subsequent matrix-product growth papers
located no resolution.

\subsection{Audit --- 2026-09-10}\label{audit-2026-09-10}

Rechecked \href{https://arxiv.org/html/2209.00449}{Varney--Morris, §7,
Question 1}. Both weak increase and weak upper regular variation remain
posed; the source explains why cancellation defeats its current
argument. Property-name and author follow-up searches found no
resolution of the two-part target.
\end{document}
